\clearpage
\fancyhead[L]{\small GAUSS AN VON LINDENAU.}
\hypertarget{native-301-lower}{}
\begin{center}GAUSS AN VON LINDENAU.
\end{center}

\begin{center}
\begin{tabular}{c}
\hline
Monatliche Correspondenz. Band XVII. S. 182--188. 1808 Februar.\\
\hline
\end{tabular}
\end{center}

\begin{flushright}
Göttingen, am 25. Januar 1808.
\end{flushright}

Entschuldigen Sie durch ein beispiellos schlechtes Wetter meinen vereitelten Wunsch, Ihnen
eine Reihe guter hiesiger Cometen-Beobachtungen mitschicken zu können. In der ganzen Zeit,
dass ich hier bin (seit dem 21. November) sind mehr nicht als 3 Abende gewesen, wo ich dem
ungünstigen Himmel einige kärgliche Beobachtungen habe abgewinnen können, die ich hier
nebst den Unterschieden meiner Elemente folgen lasse: [es sind die Beobachtungen vom
Dec. 17. 29. Jan. 4., die weiter unten Seite 298 im Abdruck des Aufsatzes vom 25. Febr.
1808 der Gött. gel. Anzeigen stehen.]

\par\indent Auf die Beobachtungen vom 29. December ist weniger zu rechnen, da das Instrument, womit
sie angestellt wurden, schlecht montirt ist. Die beiden ersten Beobachtungen vom 4. Januar
sind blos einzelne zwischen Wolken gemachte Bestimmungen, die ich aber an sich für gut
halte; die letzte ist von meinem würdigen Collegen, Hrn. Prof. \textsc{Harding}, und das Mittel aus
einigen gut harmonirenden Vergleichungen. Uebrigens vertragen die Beobachtungen des
Cometen in der letzten Zeit keine grosse Genauigkeit, da seine geringe Lichtstärke und die
Unkenntlichkeit seines Kerns nicht mehr so scharfe Beobachtung der Ein- und Austritte
erlauben. Im December-Heft der M.\,C. [S. 292 d. B.] ist übrigens bei meinen Elementen die
Neigung der Bahn vergessen, welche $63^\circ\;12'\;35''6$ ist. Ich sehe dies in dem so
eben mir communicirten Exemplare des Hrn. Hofrath \textsc{Blumenbach}.

\par\indent Für die mir gütigst mitgetheilten Mailänder Beobachtungen der \VestaAltarThreeOhOne{} bin ich Ihnen sehr verpflichtet. Mein vortrefflicher Freund, Hr. \textsc{Bessel}, hat die
Gefälligkeit gehabt, die letzten auf das schärfste mit gehöriger Rücksicht auf Refraction
zu reduciren. Hier die Resultate dieser Reductionen:


\begin{center}\fontsize{9}{11}\selectfont\setlength{\tabcolsep}{2pt}\renewcommand{\arraystretch}{1.15}\begin{tabular}{lrrrrr}
& & & & \multicolumn{2}{c}{Differenzen}\\
1807 & Mittl. Zeit & AR. & Declination & AR. & Declin.\\\hline
Sept. 5 & $7^{\text{u}}\,21^{\text{m}}\,13^{\text{s}}0$ & $211^\circ\;29'\;4''5$ & $8^\circ\;2'\;49''9$ Südl. & $-46''6$ & $-53''0$\\
8 & $7\;14\;47.5$ & $212\;49\;55.4$ & $8\;38\;18.1$ & $-54.0$ & $-59.3$\\
8 & $7\;29\;44.8$ & $212\;50\;11.5$ & $8\;38\;16.0$ & $-53.2$ & $-49.9$\\
10 & $7\;13\;14.9$ & $213\;44\;31.7$ & $9\;1\;45.8$ & $-65.5$ & $-61.4$\\
21 & $7\;4\;10.8$ & $218\;51\;35.2$ & $11\;7\;42.5$ & $-82.4$ & $-64.5$\\
26 & $6\;53\;51.8$ & $221\;15\;21.1$ & $12\;2\;59.3$ & $-99.5$ & $-66.9$\\
\end{tabular}\end{center}

\par\indent Herr Prof. \textsc{Harding} hat die Güte gehabt, diese Beobachtungen mit meinen III. Elementen zu
vergleichen. Hieraus haben sich folgende Differenzen ergeben: [die in diesem Abdruck oben
den Beobachtungen beigesetzt sind.]

\par\indent I\SourceLeaf{302}ch halte es für überflüssig, hienach jetzt noch einmal die III. Elemente zu verbessern. Die
Correctionen würden nur sehr klein sein, und die Ephemeride im September-Heft der M.\,C. wird zur
Auffindung hinreichen.

\par\indent Die neuen Elemente der \emph{Pallas}, wonach die Ephemeride im December-Heft der M.\,C. berechnet
war, sind folgende:

\begin{center}
X$^{\text{te}}$ Elemente
\end{center}


\begin{center}\fontsize{10}{12}\selectfont\setlength{\tabcolsep}{2pt}\renewcommand{\arraystretch}{1.15}\begin{tabular}{cr}
Epoche & \shortstack{Meridian von\\Göttingen}\\\hline
1802 & $143^\circ\;33'\;58''7$\\
1803 & $221\;39\;26.9$\\
1804 & $299\;57\;45.4$\\
1805 & $18\;3\;13.6$\\
1806 & $96\;8\;41.8$\\
1807 & $174\;14\;10.0$\\
1808 & $252\;32\;28.5$\\
\end{tabular}\end{center}

\begin{center}\fontsize{10}{12}\selectfont\setlength{\tabcolsep}{2pt}\renewcommand{\arraystretch}{1.15}\begin{tabular}{lr}
Tägliche tropische Bewegung .\;.\;.\;.\;.\; & $770^{.}2143$\\
Logarithm der halben grossen Achse .\;.\; & $0.4423149$\\
Excentricität .\;.\;.\;.\;.\;.\;.\;.\;.\;.\;.\; & $0.2450198$\\
Perihelium 1803 .\;.\;.\;.\;.\;.\;.\; & $121^\circ\;3'\;11''4$\\
$\AscendingNode$, 1803 .\;.\;.\;.\;.\;.\;.\;.\;.\;.\; & $172\;28\;56.9$\\
Neigung der Bahn .\;.\;.\;.\;.\;.\;.\; & $34\;37\;41.0$\\
\end{tabular}\end{center}

\par\indent Diese Elemente gründen sich blos auf die vier Oppositionen von 1803. 1804. 1805 und 1807; auf die
Beobachtungen von 1802 ist gar nicht Rücksicht genommen, da sie zur scharfen Bestimmung der
Opposition nicht brauchbar sind. Von jenen Oppositionen werden alle Längen genau dargestellt,,
die Breiten lassen sich zwar nicht genau vereinigen, doch gehen die Differenzen nicht auf $\frac12$
Minute. Übrigens wird die Pallas 1808 schwieriger zu beobachten sein als je, da sie sich gerade in
ihrem Aphelium befindet.

\Needspace{19\baselineskip}
\par\indent Vor einiger Zeit habe ich denn auch die Elemente der \emph{Juno} verbessert. Die Resultate und die
darnach berechnete Ephemeride für 1808 finden Sie beiliegend. Ich bemerke also nur noch, dass Herr
Prof. \textsc{Harding} die Güte gehabt hat, die neuen VII$^{\text{ten}}$ Elemente mit den
Beobachtungen des Herrn \textsc{Bessel} sorgfältig (jedoch mit Vernachlässigung der Parallaxe) zu
vergleichen.


\begin{center}\fontsize{9}{11}\selectfont\setlength{\tabcolsep}{2pt}\renewcommand{\arraystretch}{1.15}\begin{tabular}{lrrrrr}
1807 & \shortstack{Mittlere Zeit\\Lilienthal} & \shortstack{Berechnete\\AR.} & \shortstack{Untersch.\\R--B} & \shortstack{Berechn. südl.\\Declination} & \shortstack{Untersch.\\R--B}\\\hline
Mai 5 & $12^{\text{u}}\,46^{\text{m}}\,4^{\text{s}}$ & $240^\circ\;5'\;11''4$ & $+26''6$ & $3^\circ\;44'\;27''1$ & $+7''8$\\
7 & $12\;25\;33$ & $239\;42\;31.0$ & $+3.6$ & $3\;34\;18.0$ & $+17.6$\\
8 & $11\;30\;29$ & $239\;31\;22.3$ & $+5.6$ & $3\;29\;29.4$ & $+0.1$\\
19 & $12\;8\;26$ & $237\;17\;26.9$ & $-5.6$ & $2\;40\;34.5$ & $-3.8$\\
23 & $11\;51\;8$ & $236\;28\;5.2$ & $-9.5$ & $2\;26\;12.1$ & $+5.9$\\
24 & $11\;14\;57$ & $236\;16\;7.4$ & $-5.2$ & $2\;22\;59.9$ & $-0.2$\\
25 & $11\;54\;26$ & $236\;3\;31.0$ & $-2.6$ & $2\;19\;44.7$ & $-2.3$\\
26 & $11\;50\;58$ & $235\;51\;11.4$ & $-5.7$ & $2\;16\;41.9$ & $-16.0$\\
\end{tabular}\end{center}

\par\indent Bei der Ephemeride für die $\PallasSymbol$ finde ich den Druckfehler zu verbessern, dass die letzte Columne
nicht die Abstände der Pallas von der Erde, sondern die Logarithmen der Abstände angibt.

\par\indent Mein Werk über die Bestimmung der Bahnen der Weltkörper ist jetzt unter der Presse und wird
hoffentlich bald nach Ostern erscheinen können.

\par\indent Beigehend schicke ich Ihnen noch kleine Aberrations- und Nutations-Tafeln, welche ich schon vor
vier Jahren berechnet und mehrern Freunden mitgetheilt hatte, die dieselben bequem gefunden und
sich an ihren Gebrauch gewöhnt haben. Mir deucht, man rechnet damit unter allen Generaltafeln am
schnellsten und kürzesten. Dass dabei Sinus- und Logarithmen-Tafeln nöthig sind, kommt nicht in
Betracht, denn diese hat ja ohnehin jeder Astronom immer auf seinem Schreibtische. Das Exemplar
die\SourceLeaf{303}ser Tafeln, das ich Ihnen hier überschicke, ist nach den neuesten numerischen Angaben des Herrn
\textsc{von Zach} ganz neu berechnet. [Seite 123..128 dieses Bandes.]

\par\indent [Der geocentrische Lauf der Juno ist für Mitternacht in Göttingen vom 16. April bis 28. December
1808 mit viertägigen Zwischenräumen und zwar Gerade Aufsteigung und Abweichung auf Minuten genau,
der Logarithm des Abstandes von der Erde in vier Decimalen angegeben.]

\Needspace{18\baselineskip}
\fancyhead[L]{\small COMET VON 1807.}
\begin{center}
COMET VON 1807.
\end{center}

\begin{center}
\fontsize{10}{12}\selectfont\begin{tabular}{c}
\hline
Göttingische gelehrte Anzeigen. Stück 32. Seite 313..315. 1808 Februar 25. \\
\hline
\end{tabular}
\end{center}

Die diesen ganzen Winter hindurch herrschend gewesene ungünstige Witterung hat auf der hiesigen
Sternwarte nur wenige Beobachtungen des letzten Cometen, und meistens nur unter nicht
vortheilhaften Umständen, erlaubt. Da dieselben jetzt überall zu Ende sein werden, so stellen wir
hier alle diejenigen zusammen, die seit der Ankunft des Hrn. Prof. \textsc{Gauss} hier haben
gemacht werden können.


\begin{center}\fontsize{9}{11}\selectfont\setlength{\tabcolsep}{2pt}\renewcommand{\arraystretch}{1.15}\begin{tabular}{lrrrrr}
& & \multicolumn{2}{c}{Scheinbare} & \multicolumn{2}{c}{Unterschied}\\
& Mittlere Zeit & \shortstack{gerade\\Aufsteigung} & \shortstack{nordliche\\Abweichung} & \shortstack{gerade\\Aufsteig.} & Abweichung\\\hline
1807 Dec. 17. & $11^{\text{u}}\,24^{\text{m}}\,44^{\text{s}}$ & $305^\circ\;54'\;0''$ & $45^\circ\;4'\;30''$ & $+49$ & $-28$\\
29 & $7\;59\;21$ & $318\;57\;41$ & $46\;41\;48$ & $+6$ & $-25$\\
29 & $11\;36\;33$ & $319\;6\;35$ & $46\;43\;11$ & $+49$ & $-55$\\
1808 Jan. 4 & $7\;22\;12$ & $325\;11\;19$ & $47\;12\;46$ & $-18$ & $-21$\\
4 & $9\;13\;57$ & $325\;15\;53$ & & $-6$ & \\
4 & $17\;30\;26$ & $325\;36\;38$ & & $+18$ & \\
31 & $9\;31\;34$ & $349\;46\;43$ & $48\;9\;26$ & $-84$ & $+24$\\
\end{tabular}\end{center}

\par\indent Alle diese Beobachtungen wurden mit dem Kreismicrometer durch Vergleichung mit Sternen der
\emph{histoire céleste} angestellt, die letzte ausgenommen, wo ein bisher nicht bestimmter Stern
siebenter Grösse zur Vergleichung angewandt wurde, für dessen Stellung erst am 14. Februar eine
Beobachtung erhalten werden konnte. Diese gab


\begin{center}\fontsize{10}{12}\selectfont\setlength{\tabcolsep}{2pt}\renewcommand{\arraystretch}{1.15}\begin{tabular}{lr}
scheinbare gerade Aufsteigung des Sterns .\;.\;.\;.\; & $350^\circ\;14'\;38''$\\
scheinbare nördliche Abweichung .\;.\;.\;.\;.\;.\;.\; & $48\;4\;33$\\
\end{tabular}\end{center}

\noindent Nach dieser künftig noch zu berichtigenden Stellung ist obige Beobachtung reducirt.

\par\indent Die bereits durch die Monatliche Correspondenz bekannten parabolischen Elemente des Cometen,
welche Hr. Prof. \textsc{Gauss} schon im Anfange Novembers berechnet hatte, stimmten mit diesen
spätern Beobachtungen folgender Maassen überein [die Unterschiede sind in diesem Abdruck den
Beobachtungen oben beigefügt.]

\par\indent Es würde kaum der Mühe werth sein, nach diesen Beobachtungen jene Elemente nochmals zu verbessern,
zumal da die letzten Beobachtungen schon wegen des schwachen Lichtes des Cometen, dessen Ein- und
Austritte am 31. Januar nur mit Mühe beobachtet werden konnten, keiner vorzüglichen Genauigkeit
fähig waren. Sollte man auf andern Sternwarten nicht viel glücklicher gewesen sein, so wird sich
schwerlich über die Abweichung der Bahn dieses Cometen von der Parabel, seiner langen Sichtbar\SourceLeaf{304}keit ungeachtet, ein entscheidender Ausspruch thun lassen: gewiss ist's indess auf alle Fälle,
dass dieselbe nur sehr klein sein kann.

\par\indent Am 14. Februar war der Comet im \textsc{Herschel}'schen Telescop nur noch eben zu sehen, allein an
eine eigentliche Beobachtung war nicht mehr zu denken.
